\section[NsJail e Container ]{\texorpdfstring{NsJail e Container
\footnote{\label{fn:ref_30}https://nsjail.dev/}}{NsJail e Container }}\label{nsjail-e-container-1}

Si vuole in questo capitolo approfondire i meccansmi utilizzati da
NsJail per l'isolamento delle risorse del kernel Linux. Tutti questi
meccanismi utilizzati da NsJail sono utilizzati anche da Docker e Incus
per implementare i container e verra' visto brevemente in che modo.

\subsection[Namespaces ]{\texorpdfstring{Namespaces
\footnote{\label{fn:ref_31}https://man7.org/linux/man-pages/man7/namespaces.7.html}}{Namespaces }}\label{namespaces-2}

Un namespace e' una feature del kernel Linux che permette di
partizionare le risorse del kernel in modo che:

\begin{itemize}
\tightlist
\item
  un insieme di processi e' in grado di vedere un certo insieme di
  risorse
\item
  un'altro insieme di processi e' in grado di vedere un'altro insieme di
  risorse.
\end{itemize}

Il kernel linux supporta diverso tipi di namespace: Cgroup, IPC,
Network, Mount, PID, Time, User ed UTS.

Il \textbf{PID namespace} permette di partizionare gli ID dei processi
in modo che per esempio, su 2 PID namespace differenti possano
coesistere due processi con process ID identici. Inoltre un processo che
esegue sul pid namespace \passthrough{\lstinline!A!} non puo' vedere
quali processi sono in esecuzione sul pid namespace
\passthrough{\lstinline!B!}.

Il \textbf{Network namespace} invece offre l'isolazione delle risorse di
networking, come device, lo stack IPv4, le tabelle di routing, ecc. Un
device fisico vive all'interno di un'unico network namespace, e le
connessioni tra namespace sono gestite dai device
\passthrough{\lstinline!veth!}.\footnote{\label{fn:ref_32}https://man7.org/linux/man-pages/man7/network\_namespaces.7.html}

Banalmente, i container usano i namespace per creare un'ambiente
separato in cui istanziare i processi, il filesystem e le interfacce di
rete.

\subsection[Seccomp ]{\texorpdfstring{Seccomp
\footnote{\label{fn:ref_33}https://www.man7.org/linux/man-pages/man2/seccomp.2.html}}{Seccomp }}\label{seccomp-4}

Seccomp permette ad un processo di specificare un filtro per le chiamate
di sistema. Il filtro e' espresso come un programma
\passthrough{\lstinline!eBPF!}, che opera sul numero della syscall ed i
suoi argomenti. Fornisce un meccanismo per minimizzare l'esposizione del
kernel.

Guardiamo brevemente come si usa \passthrough{\lstinline!seccomp!},
prendendo ad esempio la specifica per il linguaggio
\passthrough{\lstinline!C!}:

\begin{itemize}
\tightlist
\item
  Si inizializza nel codice il filtro
  \passthrough{\lstinline!scmp\_filter\_ctx ctx!}
\item
  Si aggiungono le regole usando
  \passthrough{\lstinline!seccomp\_rule\_add(ctx, ...)!}
\item
  Si caricano le regole con
  \passthrough{\lstinline!seccomp\_load(ctx, ...)!}. In questo modo, il
  processo che esegue il codice sara soggetto alle limitazioni di
  \passthrough{\lstinline!seccomp!}.
\end{itemize}

I container Docker (quindi non privilegiati) usano Seccomp per limitare
le chiamate di sistema considerate pericolose, mentre Incus nei
container di sistema non applica filtri.

\subsection[cgroups ]{\texorpdfstring{cgroups
\footnote{\label{fn:ref_34}https://www.kernel.org/doc/html/latest/admin-guide/cgroup-v1/cgroups.html}}{cgroups }}\label{cgroups-5}

I Control Groups sono un meccanismo per aggregare o partizionare insiemi
di task, ed i loro figli in gruppi gerarchici.

In particolare, la documentazione del kernel Linux riporta le seguenti
definizioni:

\begin{itemize}
\tightlist
\item
  \passthrough{\lstinline!cgroup!}: un cgroup associa un insieme di task
  con un insieme di parametri per uno o piu sottosistemi.
\item
  \passthrough{\lstinline!sottosistema!}: un sottosistema e' un
  ``resource controller'' che schedula una risorsa, oppure applica
  limitazioni alla risorsa ad un certo raggruppamento di task
\end{itemize}

Dunque con un \passthrough{\lstinline!cgroup!} posso schedulare la
risorsa CPU all'interno di un gruppo definito di task. Viene utilizzato
per effettaure il tracciamento dei job, accountig e limitazione delle
risorse.

Nei container i \passthrough{\lstinline!cgroup!} sono utilizzati per
impostare un tetto massimo alle risorse utilizzabili e fare accounting
delle risorse utilizzate.
